\section{Framework 4: Uniform Control via Concentration of Measure}
\label{app:concentration_measure}

This appendix develops the concentration of measure tools required to control the fluctuations of the Yang-Mills measure uniformly in the lattice volume. This is a critical component for establishing the existence of the thermodynamic limit and the non-vanishing of the mass gap.

\subsection{Concentration on Compact Lie Groups}

The fundamental tool is the concentration of measure phenomenon on the group manifold $G = SU(N)$. Equipped with the bi-invariant metric and the Haar measure, $SU(N)$ behaves like a sphere of high dimension.

\begin{theorem}[Levy's Lemma for $SU(N)$]
Let $f: SU(N) \to \mathbb{R}$ be a Lipschitz function with constant $L$. Then for any $\epsilon > 0$:
\begin{equation}
\mu(\{ x \in SU(N) : |f(x) - \mathbb{E}[f]| \ge \epsilon \}) \le C_1 \exp\left( - \frac{C_2 N \epsilon^2}{L^2} \right)
\end{equation}
where $C_1, C_2$ are universal constants.
\end{theorem}

This implies that observables depending on many link variables are tightly concentrated around their mean values, provided they are sufficiently smooth (Lipschitz).

\subsection{Application to the Yang-Mills Action}

The lattice Yang-Mills action $S_L(U)$ is a sum of local plaquette terms. While the action itself is extensive (proportional to Volume), the \textit{fluctuations} of intensive quantities (like the energy density or the average plaquette) are suppressed by the volume.

However, for the mass gap, we need to control the fluctuations of \textit{non-local} operators (like large Wilson loops) or the correlations between distant regions.

\begin{theorem}[Uniform Concentration for Local Observables]
Let $\mathcal{O}_\Lambda$ be a sum of local operators supported on a region $\Lambda$. If the Gibbs measure satisfies a Log-Sobolev Inequality (LSI) with constant $\alpha$, then:
\begin{equation}
\text{Var}_\mu(\mathcal{O}_\Lambda) \le \frac{1}{2\alpha} \sum_{x \in \Lambda} \|\nabla_x \mathcal{O}_\Lambda\|_\infty^2
\end{equation}
\end{theorem}

\subsection{The Spectral Gap from Concentration}

The connection between concentration of measure and the spectral gap is provided by the Herbst argument.

\begin{theorem}[Gap Implies Concentration]
If the Hamiltonian $H$ has a spectral gap $\Delta > 0$, then the ground state measure $\mu_0$ satisfies a Gaussian concentration inequality:
\begin{equation}
\int e^{\lambda (f - \langle f \rangle)} d\mu_0 \le e^{\frac{\lambda^2 \|\nabla f\|_\infty^2}{2\Delta}}
\end{equation}
\end{theorem}

Conversely, establishing such concentration properties for the finite-volume measures allows us to deduce lower bounds on the spectral gap that are independent of the volume.

\subsection{Proof of Uniformity}

We use the Martingale method to prove that the concentration properties extend to the infinite volume limit.

\begin{proof}[Sketch of Uniformity Proof]
1. Consider the filtration of $\sigma$-algebras $\mathcal{F}_k$ generated by loops of size up to $k$.
2. Construct a Doob martingale for the observable $W(C)$.
3. Use the Azuma-Hoeffding inequality to bound the differences.
4. The uniform LSI (proven in Appendix \ref{app:uniform_lsi}) ensures the martingale differences decay exponentially with the distance, giving a volume-independent bound on the variance.
\end{proof}

This framework ensures that the mass gap does not close due to "entropic" fluctuations of the vacuum in the infinite volume limit.


